\documentclass[11pt,letterpaper,oneside]{article}

\usepackage[T1]{fontenc}
\usepackage[spanish,es-nodecimaldot]{babel}
\usepackage{lmodern}
\usepackage{microtype}
\usepackage{geometry}
\usepackage{xcolor}
\usepackage{graphicx}
\usepackage{booktabs}
\usepackage{tabularx}
\usepackage{longtable}
\usepackage{array}
\usepackage{enumitem}
\usepackage{listings}
\usepackage{float}
\usepackage{pdflscape}
\usepackage{caption}
\usepackage{subcaption}
\usepackage{fancyhdr}
\usepackage{titlesec}
\usepackage{xurl}
\usepackage{hyperref}
\usepackage{bookmark}

\geometry{top=2.2cm,bottom=2.2cm,left=2.35cm,right=2.35cm,headheight=15pt}
\setlength{\parindent}{0pt}
\setlength{\parskip}{0.65em}
\setlength{\emergencystretch}{2em}
\setlist{itemsep=0.25em,topsep=0.35em}
\renewcommand{\arraystretch}{1.2}

\definecolor{CenigrafNavy}{HTML}{17324D}
\definecolor{CenigrafBlue}{HTML}{2A7F9E}
\definecolor{CenigrafGold}{HTML}{B6812D}
\definecolor{CenigrafRed}{HTML}{A94F55}
\definecolor{CenigrafLight}{HTML}{EEF3F6}
\definecolor{CodeBackground}{HTML}{F4F6F8}
\definecolor{CodeComment}{HTML}{587A61}

\hypersetup{
  colorlinks=true,
  linkcolor=CenigrafNavy,
  urlcolor=CenigrafBlue,
  citecolor=CenigrafBlue,
  pdftitle={Arquitectura de estados - Proyecto Cenigraf},
  pdfauthor={Proyecto Cenigraf}
}

\titleformat{\section}{\Large\bfseries\color{CenigrafNavy}}{\thesection}{0.7em}{}
\titleformat{\subsection}{\large\bfseries\color{CenigrafBlue}}{\thesubsection}{0.7em}{}
\titleformat{\subsubsection}{\normalsize\bfseries\color{CenigrafGold}}{\thesubsubsection}{0.7em}{}

\pagestyle{fancy}
\fancyhf{}
\fancyhead[L]{\small Proyecto Cenigraf}
\fancyhead[R]{\small Arquitectura de estados en Unity}
\fancyfoot[C]{\thepage}
\renewcommand{\headrulewidth}{0.4pt}

\lstdefinestyle{cenigraf}{
  language=[Sharp]C,
  backgroundcolor=\color{CodeBackground},
  basicstyle=\ttfamily\small,
  keywordstyle=\color{CenigrafBlue}\bfseries,
  commentstyle=\color{CodeComment},
  stringstyle=\color{CenigrafRed},
  numbers=left,
  numberstyle=\scriptsize\color{gray},
  stepnumber=1,
  numbersep=8pt,
  frame=single,
  rulecolor=\color{CenigrafLight},
  breaklines=true,
  showstringspaces=false,
  tabsize=4,
  columns=fullflexible
}
\lstset{style=cenigraf}

\newcommand{\implemented}{\textbf{Implementado}}
\newcommand{\planned}{\textbf{Planeado}}
\newcommand{\code}[1]{\nolinkurl{#1}}
\newcommand{\figurepath}{images}

\begin{document}

\hypersetup{pageanchor=false}
\begin{titlepage}
  \pagecolor{CenigrafNavy}
  \color{white}
  \vspace*{2.1cm}
  {\Huge\bfseries Arquitectura de estados\par}
  \vspace{0.35cm}
  {\LARGE Aplicación en Unity para el Proyecto Cenigraf\par}
  \vspace{1cm}
  {\color{CenigrafGold}\rule{\textwidth}{2pt}}
  \vspace{1cm}
  {\large Diseño, flujo de datos, contratos y ruta de implementación\par}
  \vfill
  \begin{tabular}{@{}ll@{}}
    \textbf{Motor:} & Unity 6 (6000.3.18f1) \\
    \textbf{Lenguaje:} & C\# \\
    \textbf{Documento:} & Versión 1.0 \\
    \textbf{Fecha:} & Agosto de 2026
  \end{tabular}
  \vspace{1.4cm}
\end{titlepage}
\hypersetup{pageanchor=true}
\nopagecolor
\color{black}

\pagenumbering{roman}
\tableofcontents
\clearpage
\listoffigures
\clearpage
\pagenumbering{arabic}

\section{Propósito y alcance}

Este documento define la arquitectura de estados del Proyecto Cenigraf y su aplicación en Unity. Su objetivo es establecer un lenguaje común entre diseño, programación y pruebas para representar el mundo, el jugador y los objetos decorativos mediante datos explícitos y transiciones controladas.

La propuesta toma como fuente los diagramas entregados y contrasta sus conceptos con el código C\# disponible en el repositorio. Esta distinción es importante: los diagramas describen la visión objetivo, mientras que las clases del proyecto muestran qué partes ya funcionan en Unity.

\subsection{Alcance funcional}

La documentación cubre:

\begin{itemize}
  \item composición del estado principal y de los estados \code{World}, \code{Player} y \code{ObjDecoration};
  \item responsabilidades de \code{StateData}, \code{Action}, \code{Payload}, \code{Reducer}, \code{Store} y adaptadores de runtime;
  \item flujo de movimiento y reconciliación con \code{Rigidbody2D};
  \item correspondencia entre los campos planeados y tipos apropiados de Unity/C\#;
  \item estrategia incremental para completar los estados aún no implementados;
  \item reglas de validación, persistencia y pruebas.
\end{itemize}

No se describe aquí la lógica de interfaz, autenticación o publicación de eventos, excepto cuando consumen o producen cambios de estado.

\subsection{Convenciones de estado de avance}

\begin{table}[H]
\centering
\caption{Estados usados para describir el avance técnico}
\begin{tabularx}{\textwidth}{>{\bfseries}p{2.6cm}X}
\toprule
Etiqueta & Significado \\
\midrule
Implementado & Existe una representación funcional en el código actual. \\
Parcial & Existe la infraestructura base, pero el modelo no contiene todavía todos los campos del diseño. \\
Planeado & Está definido en los diagramas o en esta especificación, pero requiere implementación. \\
\bottomrule
\end{tabularx}
\end{table}

\subsection{Vista general de la composición}

El estado principal se concibe como un contenedor de subestados independientes. Esta división evita que una actualización del jugador tenga que reconstruir o conocer todos los detalles del mundo, y permite evolucionar cada dominio con acciones y reductores propios.

\begin{figure}[H]
  \centering
  \includegraphics[width=\textwidth]{\figurepath/composicion-estados.png}
  \caption{Composición planeada del estado principal, el mundo, el jugador y los objetos decorativos.}
  \label{fig:composicion-estados}
\end{figure}

La Figura~\ref{fig:composicion-estados} funciona como mapa conceptual, no como definición literal de tipos C\#. En las secciones siguientes se normalizan nombres, tipos y responsabilidades para que el diseño sea compatible con serialización, física 2D y convenciones de Unity.

\section{Patrón arquitectónico}

La arquitectura adopta un flujo unidireccional inspirado en Redux, ajustado al ciclo de vida de Unity. Los datos actuales residen en un \emph{store}; cualquier intención de cambio se expresa como una acción; un reductor calcula el siguiente estado; finalmente, un adaptador aplica al runtime los efectos necesarios.

\subsection{Componentes y responsabilidades}

\begin{table}[H]
\centering
\caption{Responsabilidades de los componentes del patrón}
\begin{tabularx}{\textwidth}{>{\bfseries}p{3.1cm}X X}
\toprule
Componente & Responsabilidad & Restricción principal \\
\midrule
StateData & Contener una instantánea serializable del dominio. & No controla objetos de escena ni ejecuta efectos. \\
Action & Nombrar una intención: mover, cambiar rol, configurar colisión. & Debe ser explícita y tener un tipo inequívoco. \\
Payload & Transportar los valores necesarios para la acción. & Solo datos; sin acceso al store. \\
Reducer & Calcular el siguiente estado desde el estado actual y la acción. & Debe ser determinista y no modificar la escena. \\
Store & Custodiar el estado, ejecutar el reductor y publicar cambios. & Es el único punto de escritura del estado. \\
Runtime Applier & Traducir el estado a componentes Unity. & Escucha al store; no decide reglas de dominio. \\
Controlador & Lee entrada o eventos y despacha acciones. & No escribe directamente el StateData. \\
\bottomrule
\end{tabularx}
\end{table}

\subsection{Flujo unidireccional}

\begin{figure}[H]
  \centering
  \includegraphics[width=0.88\textwidth]{\figurepath/flujo-datos.png}
  \caption{Flujo conceptual desde el estado actual hasta la actualización del store y la reconciliación física.}
  \label{fig:flujo-datos}
\end{figure}

El flujo esperado es:

\begin{enumerate}
  \item un controlador captura una intención del usuario o un evento del juego;
  \item calcula o solicita un valor permitido, por ejemplo una posición sin penetrar colisiones;
  \item crea una \code{Action} con su \code{Payload};
  \item el \code{Store.Dispatch} entrega la acción al \code{Reducer};
  \item el reductor valida reglas simples y devuelve una nueva instantánea;
  \item si el estado cambió, el store publica \code{StateChanged};
  \item los adaptadores actualizan \code{Rigidbody2D}, \code{BoxCollider2D}, renderizadores u otros componentes;
  \item cuando Unity altera el valor final por física, se despacha una acción de reconciliación.
\end{enumerate}

\subsection{Propiedades del diseño}

\begin{itemize}
  \item \textbf{Trazabilidad:} cada mutación se identifica mediante un tipo de acción.
  \item \textbf{Repetibilidad:} el mismo estado y la misma acción producen el mismo resultado.
  \item \textbf{Desacoplamiento:} la lógica de datos no necesita referencias a objetos de escena.
  \item \textbf{Pruebas simples:} reductores y constructores se prueban sin cargar una escena.
  \item \textbf{Integración segura:} los efectos de Unity quedan concentrados en controladores y adaptadores.
\end{itemize}

\subsection{Inmutabilidad práctica}

En el código actual, \code{PlayerStateData} y \code{GameSessionStateData} son estructuras serializables con campos privados y propiedades de solo lectura. Los métodos \code{With...} construyen una nueva instantánea en vez de exponer setters. Este enfoque debe mantenerse para los estados planeados.

\begin{lstlisting}[caption={Forma recomendada de producir una nueva instantánea},label={lst:with-pattern}]
public WorldStateData WithRoom(string value) =>
    new WorldStateData(id, name, value, mainTheme, dimensions);
\end{lstlisting}

El uso de métodos \code{set.X} y \code{get.X} en los diagramas representa la intención funcional. En C\# se recomienda expresarla mediante fábricas de acciones, propiedades de solo lectura y métodos \code{With...}; no mediante setters públicos que permitan saltarse el store.

\section{Composición de los estados}

\subsection{Estado principal}

El diseño plantea un agregado raíz con los subestados \code{world}, \code{player} y una colección de \code{objDecoration}. En el proyecto actual no existe todavía un único \code{PrincipalStateData}; hay stores separados para sesión, jugador y cámara. Esta separación ya ofrece independencia por dominio y puede conservarse, incorporando un coordinador solamente si aparecen operaciones que necesiten una instantánea global consistente.

\begin{table}[H]
\centering
\caption{Mapa de dominios y avance actual}
\begin{tabularx}{\textwidth}{>{\bfseries}p{3.3cm}p{2.5cm}X}
\toprule
Dominio & Estado & Evidencia en el proyecto \\
\midrule
Sesión & Implementado & \code{GameSessionStateData}, acciones, reductor y store persistente entre escenas. \\
Jugador & Parcial & Posición, rol, ordenamiento Y y configuración del collider; faltan identidad, apariencia e inventario del diseño. \\
Cámara & Implementado & Posición, objetivo, retardo, suavizado, ajuste a píxel y zoom. \\
Mundo & Planeado & Existe lógica de renderizado por Y, pero no un \code{WorldStateData} ni su store. \\
ObjDecoration & Planeado & Los diagramas definen sus datos; no hay un store/reductor específico localizado. \\
Estado principal & Planeado/opcional & Puede materializarse como fachada de stores o como agregado serializable para guardado. \\
\bottomrule
\end{tabularx}
\end{table}

\subsection{WorldStateData}

El estado del mundo identifica el escenario activo y sus dimensiones lógicas. No debe almacenar referencias directas a escenas o GameObjects si se pretende persistir o sincronizar.

\begin{table}[H]
\centering
\caption{Campos planeados para el estado del mundo}
\begin{tabularx}{\textwidth}{>{\ttfamily}p{2.6cm}>{\ttfamily}p{3.2cm}X}
\toprule
\normalfont\bfseries Campo & \normalfont\bfseries Tipo recomendado & \normalfont\bfseries Descripción \\
\midrule
id & string & Identificador estable del mundo; se prefiere texto sobre un número dependiente de base de datos. \\
name & string & Nombre visible o administrativo. \\
room & string & Sala, zona o escena lógica activa. \\
mainTheme & string & Identificador del tema visual principal. \\
dimensions & Vector2Int & Ancho y alto del mapa en celdas o unidades lógicas. \\
\bottomrule
\end{tabularx}
\end{table}

Las acciones mínimas serían \code{SetWorld}, \code{SetRoom}, \code{SetMainTheme} y \code{SetDimensions}. Para cargar un mundo completo conviene una sola acción con un payload compuesto, evitando que los consumidores observen configuraciones intermedias incoherentes.

\begin{figure}[H]
  \centering
  \includegraphics[width=0.9\textwidth]{\figurepath/estructura-world.png}
  \caption{Estructura de datos y operaciones propuestas para el estado World.}
  \label{fig:world}
\end{figure}

\subsection{PlayerStateData}

El jugador tiene una implementación parcial que ya controla los valores con efecto inmediato en física y renderizado. El modelo objetivo agrega identidad, progresión, apariencia e inventario.

\begin{longtable}{>{\ttfamily}p{2.7cm}>{\ttfamily}p{3cm}p{6.6cm}}
\caption{Campos del jugador: diseño y normalización recomendada}\label{tab:player-fields}\\
\toprule
\normalfont\bfseries Campo & \normalfont\bfseries Tipo recomendado & \normalfont\bfseries Propósito y estado \\
\midrule
\endfirsthead
\toprule
\normalfont\bfseries Campo & \normalfont\bfseries Tipo recomendado & \normalfont\bfseries Propósito y estado \\
\midrule
\endhead
id & string & Identidad estable. Planeado. \\
nickname & string & Nombre público. Planeado. \\
email & string & Dato de cuenta; no debería viajar en estados públicos de red. Planeado. \\
level & int & Progresión; validar un mínimo de cero o uno. Planeado. \\
frameAnimation & int & Índice de cuadro si realmente es dato persistente; de lo contrario pertenece al runtime. Planeado. \\
skinId & string & Identificador de apariencia. El diagrama usa \code{rana/skin}; se recomienda un nombre único. Planeado. \\
roleId & string & Rol elegido en la sesión. Implementado. \\
position & Vector2 & Posición lógica del cuerpo 2D. Implementado. \\
colliderSize & Vector2 & Dimensiones del \code{BoxCollider2D}. Implementado con mínimo de 0.01. \\
colliderOffset & Vector2 & Desplazamiento del collider respecto al transform. Implementado. \\
ySortEnabled & bool & Activa el ordenamiento visual por coordenada Y. Implementado. \\
dimensions & PlayerSpatialData & Agrupación opcional para sprites, colliders e iconos; planeada y por definir. \\
inventoryEnabled & bool & Permite activar el subsistema de inventario. Planeado. \\
inventory & InventoryStateData & Lista o agregado de ítems identificables. Planeado. \\
\bottomrule
\end{longtable}

Los bloques \code{dimension}, \code{position} y \code{offset} del diagrama mezclan matrices, valores escalares e información para varios sistemas. En Unity conviene separar semánticas: \code{Vector2} para medidas/posiciones, arreglos tipados para datos por capa y estructuras específicas cuando un bloque crece. Una matriz \code{[X,Y]} solo es apropiada si representa una cuadrícula real.

\begin{figure}[H]
  \centering
  \includegraphics[height=0.82\textheight]{\figurepath/estructura-player.png}
  \caption{Estructura planeada del jugador y catálogo conceptual de operaciones.}
  \label{fig:player}
\end{figure}

\subsection{ObjDecorationStateData}

Un objeto decorativo describe elementos del escenario que pueden tener representación visual, colisión, interacción o inventario. Para mantener el sistema escalable, cada instancia necesita identidad propia y los objetos repetidos deben almacenarse en una colección indexada por \code{id}.

\begin{table}[H]
\centering
\caption{Campos planeados para objetos decorativos}
\begin{tabularx}{\textwidth}{>{\ttfamily}p{2.7cm}>{\ttfamily}p{3.1cm}X}
\toprule
\normalfont\bfseries Campo & \normalfont\bfseries Tipo recomendado & \normalfont\bfseries Descripción \\
\midrule
id & string & Identificador único y estable. \\
imageId & string & Clave de sprite, Addressable o recurso; evita serializar un \code{Sprite} dentro del estado puro. \\
themeId & string & Tema o variante visual. \\
position & Vector2 & Posición lógica en el mundo. \\
dimensions & Vector2 & Tamaño visual o de celda. \\
colliderSize & Vector2 & Tamaño del collider, separado del sprite. \\
colliderOffset & Vector2 & Desplazamiento local del collider. \\
lineIndex & int & Índice de línea/capa si el mapa lo requiere; para Y-sort puede derivarse de la posición. \\
areaInfo & AreaInfoState[] & Regiones informativas. \\
interactions & InteractionState[] & Interacciones disponibles. \\
icons & string[] & Identificadores de iconos asociados. \\
inventoryEnabled & bool & Controla si el objeto puede contener ítems. \\
inventory & InventoryStateData & Contenido del objeto; preferiblemente IDs y cantidades. \\
\bottomrule
\end{tabularx}
\end{table}

\begin{figure}[H]
  \centering
  \includegraphics[height=0.82\textheight]{\figurepath/estructura-obj-decoration.png}
  \caption{Estructura planeada de ObjDecoration y operaciones de actualización.}
  \label{fig:obj-decoration}
\end{figure}

Las operaciones repetidas que aparecen en el diagrama deben consolidarse en un catálogo único de acciones. Además, \code{enableInventory} es booleano y no un objeto; la acción correspondiente debería aceptar \code{bool}. Los nombres \code{imagebg} y \code{theme} deberían normalizarse a \code{imageId} y \code{themeId} para indicar que son referencias lógicas.

\section{Aplicación del flujo en Unity}

\subsection{Inicialización}

La inicialización debe tomar como fuente los componentes existentes de la escena y producir una instantánea válida antes de procesar entrada. El proyecto ya aplica este criterio en \code{PlayerStateRuntimeApplier.Awake}: obtiene \code{Rigidbody2D}, \code{BoxCollider2D}, \code{WorldYSort} y el rol de la sesión, inicializa el store y aplica el estado.

El orden de ejecución relevante es:

\begin{enumerate}
  \item \code{GameSessionStore} se crea o despierta con orden $-1000$ y sobrevive cambios de escena.
  \item Los objetos de la nueva escena ejecutan \code{Awake}.
  \item \code{PlayerStateRuntimeApplier}, con orden $-100$, construye el estado inicial del jugador.
  \item Los controladores comienzan a leer entrada en \code{Update} y a mover en \code{FixedUpdate}.
\end{enumerate}

\subsection{Movimiento del jugador}

El movimiento implementado usa la física como restricción y el store como fuente lógica de posición:

\begin{enumerate}
  \item \code{Update} lee la acción \code{Player/Move} y limita su magnitud.
  \item \code{FixedUpdate} reconcilia un movimiento físico previo si Unity produjo una posición distinta.
  \item Se calcula el desplazamiento solicitado con velocidad y \code{fixedDeltaTime}.
  \item \code{Physics2D.BoxCast} resuelve cada eje de forma independiente para permitir deslizamiento contra paredes.
  \item Se despacha \code{PlayerAction.SetPosition(allowedPosition)}.
  \item El store reduce y publica el nuevo estado.
  \item \code{Rigidbody2D.MovePosition} solicita el movimiento físico.
  \item En el siguiente paso fijo se usa \code{ReconcilePosition} cuando la posición física final difiere de la lógica.
\end{enumerate}

Esta reconciliación evita que el estado quede desfasado cuando otro cuerpo, un contacto o el propio motor ajusta la posición solicitada.

\begin{lstlisting}[caption={Secuencia esencial implementada para el movimiento},label={lst:movement}]
Vector2 allowedPosition = ResolveAllowedPosition(
    store.State.Position,
    requestedDisplacement);

store.Dispatch(PlayerAction.SetPosition(allowedPosition));
body.MovePosition(store.State.Position);
\end{lstlisting}

\subsection{Aplicación a componentes de runtime}

El \code{PlayerStateRuntimeApplier} escucha \code{StateChanged} y refleja en Unity el ordenamiento Y y los datos del collider. La posición se aplica desde el controlador de movimiento, porque debe sincronizarse con \code{FixedUpdate}. Esta división es válida si se documenta la propiedad de cada efecto:

\begin{table}[H]
\centering
\caption{Propietario recomendado de cada efecto}
\begin{tabularx}{\textwidth}{>{\bfseries}p{3.3cm}p{4.2cm}X}
\toprule
Dato & Aplicador & Motivo \\
\midrule
Posición de Rigidbody2D & Controlador de movimiento & Debe respetar el paso fijo y las consultas de colisión. \\
Tamaño/offset de collider & Runtime applier & Es una proyección directa del estado. \\
Y-sort & Runtime applier / WorldYSort & Es un efecto visual derivado. \\
Skin y animación & Adaptador visual del jugador & Aísla Animator, SpriteRenderer y carga de assets. \\
Decoraciones & World/Decoration runtime registry & Crea, recicla o actualiza vistas según IDs. \\
Inventario & Adaptador de inventario/UI & La vista reacciona; las reglas permanecen en el dominio. \\
\bottomrule
\end{tabularx}
\end{table}

\subsection{Eventos y prevención de ciclos}

Un consumidor de \code{StateChanged} puede aplicar efectos, pero no debería volver a despachar la misma acción sin una condición de salida. Para evitar ciclos:

\begin{itemize}
  \item el store ignora transiciones cuyo resultado es igual al estado actual;
  \item las acciones de reconciliación usan tolerancia numérica;
  \item los aplicadores no leen entrada ni inventan valores de dominio;
  \item las suscripciones se realizan en \code{OnEnable} y se eliminan en \code{OnDisable}.
\end{itemize}

\subsection{Persistencia y cambio de escena}

El estado de sesión ya usa \code{DontDestroyOnLoad}. Para el estado global planeado se recomiendan dos niveles:

\begin{itemize}
  \item \textbf{Estado de sesión:} selección de rol, sala activa y datos necesarios mientras la aplicación está abierta.
  \item \textbf{Estado persistente:} progreso, inventarios y configuraciones que se serializan a JSON o a un repositorio.
\end{itemize}

No se deben serializar referencias directas a \code{GameObject}, \code{Transform}, \code{Collider2D} o \code{Sprite}. En su lugar se guardan IDs y datos simples; al cargar una escena, un registro de runtime resuelve esos IDs a componentes o assets.

\section{Ruta de implementación propuesta}

La evolución debe ser incremental para conservar el comportamiento actual. Cada etapa puede integrarse y probarse de manera independiente.

\subsection{Etapa 1: consolidar contratos y nombres}

\begin{enumerate}
  \item Definir una convención: tipos en \code{PascalCase}, campos serializados en \code{camelCase} y acciones con verbo explícito.
  \item Resolver términos ambiguos del diagrama: \code{skinId}, \code{imageId}, \code{themeId}, \code{dimensions}, \code{colliderSize} y \code{colliderOffset}.
  \item Definir IDs como \code{string} si pueden provenir de servicios o bases de datos heterogéneas.
  \item Crear contratos de inventario, interacción y área antes de agregarlos a los estados principales.
\end{enumerate}

\subsection{Etapa 2: completar PlayerState}

Agregar identidad y apariencia al modelo existente sin reemplazar la ruta de movimiento. Cada dato nuevo requiere, como mínimo:

\begin{itemize}
  \item campo privado serializado y propiedad de solo lectura;
  \item parámetro de constructor con valor seguro;
  \item método \code{With...};
  \item tipo de acción, payload y fábrica de acción;
  \item caso en el reductor;
  \item comparación en \code{StatesAreEqual};
  \item prueba del reductor y, si produce un efecto visual, prueba del adaptador.
\end{itemize}

El correo electrónico debe permanecer en el estado de cuenta o perfil si no es necesario para la simulación del jugador. Separarlo reduce exposición accidental y evita replicar datos privados a otros clientes.

\subsection{Etapa 3: implementar WorldState}

Crear el módulo \code{World/State} con \code{WorldStateData}, \code{WorldAction}, \code{WorldStateReducer} y \code{WorldStateStore}. Su primera versión puede manejar mundo, sala, tema y dimensiones. La carga completa debe ser atómica:

\begin{lstlisting}[caption={Ejemplo de acción compuesta para cargar un mundo}]
public static WorldAction LoadWorld(WorldDefinition definition) =>
    new WorldAction(
        WorldActionType.LoadWorld,
        new WorldPayload(
            definition.Id,
            definition.DisplayName,
            definition.RoomId,
            definition.ThemeId,
            definition.Dimensions));
\end{lstlisting}

\subsection{Etapa 4: implementar decoraciones como colección}

Un único \code{ObjDecorationStateData} describe una instancia. El store de mundo debería custodiar una colección o un diccionario por ID y ofrecer acciones como:

\begin{itemize}
  \item \code{AddDecoration(payload)};
  \item \code{RemoveDecoration(id)};
  \item \code{MoveDecoration(id, position)};
  \item \code{SetDecorationTheme(id, themeId)};
  \item \code{SetDecorationInventory(id, inventory)}.
\end{itemize}

Para actualizaciones frecuentes, un diccionario facilita buscar por ID. Para serialización con el \code{JsonUtility} tradicional de Unity, una lista es más directa; puede mantenerse una lista serializable y construir un índice de runtime no serializado.

\subsection{Etapa 5: fachada de estado principal}

Solo cuando un caso real lo requiera, crear una fachada que exponga una instantánea coherente de sesión, mundo, jugador y decoraciones. No es obligatorio fusionar todos los stores en un único \code{MonoBehaviour}. Una fachada de solo lectura puede componerlos:

\begin{lstlisting}[caption={Fachada conceptual de solo lectura}]
public readonly struct PrincipalStateSnapshot
{
    public GameSessionStateData Session { get; }
    public WorldStateData World { get; }
    public PlayerStateData Player { get; }
    public IReadOnlyList<ObjDecorationStateData> Decorations { get; }
}
\end{lstlisting}

Esta opción conserva la independencia de cada ciclo de actualización y permite producir guardados, telemetría o depuración global sin habilitar mutaciones externas.

\subsection{Mapa completo de referencia}

La Figura~\ref{fig:mapa-completo} reúne los estados y los métodos previstos. Debido a su relación de aspecto, se presenta en una página horizontal para conservar legibilidad.

\begin{landscape}
\begin{figure}[p]
  \centering
  \includegraphics[width=0.98\linewidth,height=0.86\textheight,keepaspectratio]{\figurepath/mapa-estados-metodos.png}
  \caption{Mapa consolidado de estados, propiedades y operaciones planeadas.}
  \label{fig:mapa-completo}
\end{figure}
\end{landscape}

Los métodos del mapa deben interpretarse como inventario de capacidades. En la implementación, los getters son propiedades de solo lectura y los setters se convierten en acciones; no se recomienda copiar literalmente pares \code{set/get} para cada campo.

\section{Validación, pruebas y observabilidad}

\subsection{Invariantes por dominio}

Los reductores deben proteger invariantes baratas y deterministas. Las reglas que necesitan acceso a escena, red o base de datos pertenecen a servicios previos al despacho.

\begin{table}[H]
\centering
\caption{Invariantes mínimas recomendadas}
\begin{tabularx}{\textwidth}{>{\bfseries}p{3cm}X X}
\toprule
Dominio & Invariante & Comportamiento esperado \\
\midrule
Player & Cada dimensión del collider es al menos 0.01. & Limitar en el reductor; ya implementado. \\
Player & El nivel no es negativo. & Limitar o rechazar la acción según la regla de negocio. \\
Player & La posición contiene valores finitos. & Rechazar NaN e infinito antes de actualizar. \\
World & Ancho y alto son positivos. & Rechazar mundos sin área válida. \\
Decoration & El ID no está vacío y no se duplica. & Validar al agregar la instancia. \\
Decoration & Inventario solo se usa cuando está habilitado. & Vaciar, ignorar o conservar según una política documentada. \\
Todos & IDs y textos nulos se normalizan. & Convertir a cadena vacía o rechazar según obligatoriedad. \\
\bottomrule
\end{tabularx}
\end{table}

\subsection{Pirámide de pruebas}

\subsubsection{Pruebas unitarias de reductores}

Son la base del sistema. Cada tipo de acción debe probar estado inicial, payload, resultado y preservación de campos no relacionados. También deben cubrir valores límite y acciones desconocidas.

\subsubsection{Pruebas de stores}

Verifican que \code{StateChanged} se emite una sola vez cuando hay cambio y no se emite para una transición equivalente. Deben comprobar que el evento entrega la acción que causó la transición.

\subsubsection{Pruebas de integración en Play Mode}

Cubren la conexión con \code{Rigidbody2D}, colliders, \code{WorldYSort}, animación y cambios de escena. Casos esenciales:

\begin{itemize}
  \item el jugador no atraviesa obstáculos y puede deslizarse por sus bordes;
  \item la reconciliación conserva alineados store y cuerpo físico;
  \item cambiar el tamaño del collider actualiza el componente y respeta el mínimo;
  \item cambiar de escena preserva la selección de rol;
  \item crear o quitar una decoración sincroniza el registro visual sin duplicados.
\end{itemize}

\subsection{Registro de transiciones}

Durante desarrollo, un observador opcional puede registrar \code{timestamp}, store, tipo de acción y resumen del estado resultante. No debe incluir correo, tokens ni inventarios sensibles. Para estados grandes se registra un resumen o diferencia, no la instantánea completa en cada frame.

\subsection{Rendimiento}

El patrón no implica despachar todos los valores visuales continuamente. Recomendaciones:

\begin{itemize}
  \item no guardar en el estado global cuadros de animación que cambian cada frame, salvo que sean necesarios para red o repetición;
  \item evitar reconstruir colecciones completas para una sola decoración en escenarios masivos; usar actualización por ID y estructuras adecuadas;
  \item no publicar eventos si \code{StatesAreEqual} confirma que no hubo cambio;
  \item mantener la física en \code{FixedUpdate} y la lectura de entrada en \code{Update};
  \item perfilar antes de optimizar o reemplazar estructuras claras por soluciones complejas.
\end{itemize}

\subsection{Criterios de aceptación de la arquitectura}

La implementación objetivo se considera completa cuando:

\begin{enumerate}
  \item todo estado se modifica únicamente mediante acciones y reductores;
  \item los estados planeados tienen tipos C\# concretos y serializables;
  \item cada campo tiene propietario, significado, unidad y valor inicial documentados;
  \item los efectos de Unity están fuera de los reductores;
  \item las transiciones críticas cuentan con pruebas unitarias y de integración;
  \item el guardado y la carga reconstruyen una instantánea equivalente;
  \item los IDs resuelven assets y objetos sin serializar referencias de escena.
\end{enumerate}

\section{Anexos}

\subsection{Correspondencia con el código actual}

Las clases citadas se encuentran bajo \code{Assets/\_SocialProyectoCenigraf/Scripts}:

\begin{longtable}{>{\ttfamily}p{5.8cm}p{7cm}}
\caption{Clases implementadas relevantes para la arquitectura}\label{tab:current-code}\\
\toprule
\normalfont\bfseries Clase & \normalfont\bfseries Responsabilidad observada \\
\midrule
\endfirsthead
\toprule
\normalfont\bfseries Clase & \normalfont\bfseries Responsabilidad observada \\
\midrule
\endhead
GameSessionStateData & Conserva el ID del rol seleccionado. \\
GameSessionAction & Declara selección y limpieza de rol con payload tipado. \\
GameSessionReducer & Produce la nueva sesión de forma determinista. \\
GameSessionStore & Store singleton persistente entre escenas. \\
PlayerStateData & Conserva posición, rol, Y-sort, tamaño y offset del collider. \\
PlayerAction & Declara posición, traslación, reconciliación, rol, Y-sort y collider. \\
PlayerStateReducer & Calcula el siguiente estado y limita dimensiones de collider. \\
PlayerStateStore & Custodia el estado del jugador y publica cambios reales. \\
PlayerStateRuntimeApplier & Inicializa y aplica collider y Y-sort al runtime. \\
PlayerMovementController & Lee input, resuelve colisiones, despacha posición y reconcilia física. \\
CameraStateData & Conserva estado lógico de posición, seguimiento y zoom de cámara. \\
CameraAction/Reducer/Store & Implementan el mismo flujo unidireccional para cámara. \\
WorldYSort & Deriva \code{sortingOrder} desde la coordenada Y. \\
\bottomrule
\end{longtable}

\subsection{Glosario}

\begin{description}[style=nextline,leftmargin=3.1cm,labelwidth=2.8cm]
  \item[Acción] Mensaje inmutable que identifica una intención de cambio.
  \item[Payload] Datos específicos transportados por una acción.
  \item[Reductor] Función determinista que calcula el siguiente estado.
  \item[Store] Componente propietario de una instantánea y del método de despacho.
  \item[Instantánea] Valor completo del estado en un momento determinado.
  \item[Reconciliación] Corrección del estado lógico con el resultado definitivo del runtime o servidor.
  \item[Runtime applier] Adaptador que traduce datos del estado a componentes Unity.
  \item[Y-sort] Ordenamiento visual basado en la posición vertical para simular profundidad 2D.
\end{description}

\subsection{Decisiones pendientes}

Antes de implementar todos los diagramas, el equipo debe resolver:

\begin{itemize}
  \item si el ID canónico es texto o número en los servicios externos;
  \item qué representa exactamente cada matriz \code{[X,Y]} del diseño;
  \item si \code{frameAnimation} es persistente, replicado o exclusivamente visual;
  \item si el inventario vive en jugador/decoración o en un store normalizado por propietario;
  \item qué datos del jugador son privados y cuáles pueden compartirse con otros clientes;
  \item si la sala equivale a una escena Unity, una zona dentro de escena o una sesión de red;
  \item si el estado principal será un store único o una fachada de stores por dominio.
\end{itemize}

\subsection{Registro de fuentes visuales}

Los seis diagramas suministrados se copiaron a \code{\_Doc/images} con nombres normalizados. Todas las figuras forman parte del proyecto documental y usan rutas relativas, por lo que no dependen de la carpeta de descargas del autor.

\begin{table}[H]
\centering
\caption{Archivos visuales incorporados}
\begin{tabularx}{\textwidth}{>{\ttfamily}p{6.4cm}X}
\toprule
\normalfont\bfseries Archivo & \normalfont\bfseries Uso \\
\midrule
composicion-estados.png & Vista general de subestados y campos. \\
flujo-datos.png & Flujo unidireccional y ejemplo de movimiento. \\
mapa-estados-metodos.png & Mapa consolidado de propiedades y operaciones. \\
estructura-world.png & Detalle de World y su payload conceptual. \\
estructura-player.png & Detalle de Player y sus operaciones. \\
estructura-obj-decoration.png & Detalle de ObjDecoration y sus operaciones. \\
\bottomrule
\end{tabularx}
\end{table}


\end{document}
